\documentclass[11pt]{article}
\usepackage[T2A]{fontenc}
\usepackage[utf8]{inputenc}
\usepackage[ukrainian,english]{babel}
\usepackage[margin=1in]{geometry}
\usepackage{amsmath,amssymb}
\usepackage{enumitem}
\usepackage{hyperref}
\usepackage{tikz}
\usepackage{microtype}
\usepackage{longtable}
\usepackage{booktabs}

\setlength{\parindent}{0pt}
\setlength{\parskip}{0.6\baselineskip}
\setlength{\emergencystretch}{3em}
\hbadness=10000
\sloppy

\title{MCM-Gen та MCM-Bench v1.0: архітектура, генерація та валідація процедурного бенчмарку для символьної регресії}
\author{Блажків Владислав Олегович\\аспірант групи phd-f3-25 «Комп'ютерні науки»}
\date{}

\begin{document}
\selectlanguage{ukrainian}

\maketitle

\section*{Анотація}

Статична природа поширених бенчмарків для символьної регресії (зокрема Feynman Dataset і SRBench) створює суттєвий методологічний ризик в оцінюванні великих мовних моделей (LLM) через можливе забруднення даних (data contamination). У попередній роботі було запропоновано теоретичну Багатовимірну матрицю складності (Multidimensional Complexity Matrix, MCM), яка класифікує вирази за трьома осями: структурною (A), семантичною (B) та топологічною (C). У цій статті подано практичну реалізацію концепції у вигляді фреймворку процедурної генерації MCM-Gen та публічного набору задач MCM-Bench v1.0. Описано алгоритмічний конвеєр генерації: шаблонний синтез дерев виразів, гібридну символьно-чисельну топологічну валідацію з керуванням тайм-аутами та адаптивне семплювання області визначення. MCM-Bench v1.0 містить 818 унікальних позарозподільних задач, що покривають 64 комірки простору $\langle A,B,C\rangle$. За результатами зовнішнього аудиту відповідності частка коректних маркувань становить 85\% для семантичної осі та 89\% для топологічної осі, що вказує на практичну придатність бенчмарку в межах заданого протоколу перевірки.

\textbf{Ключові слова:} символьна регресія, великі мовні моделі, процедурна генерація, бенчмаркінг, топологічна валідація, MCM-Gen, MCM-Bench.

\section{Вступ}

Оцінювання LLM у задачах символьної регресії перебуває в точці методологічної перебудови. Високі результати на статичних наборах даних дедалі важче інтерпретувати як прояв реального математичного міркування, оскільки відкриті колекції формул можуть потрапляти до передтренувальних корпусів моделей. У нашій попередній теоретичній роботі було сформульовано підхід MCM, який відокремлює складність задач за трьома ортогональними вимірами й тим самим створює основу для процедурного, а не архівного, тестування.

Ключова ідея MCM --- інверсія класичного сценарію <<дані $\rightarrow$ формула>> у сценарій <<формула $\rightarrow$ дані>> (Reverse-Flow Generation). Спочатку синтезується аналітичний вираз із контрольованими властивостями, далі незалежно верифікуються його структурні, семантичні та топологічні ознаки, і лише після цього генеруються числові спостереження. Такий підхід спрямований на зменшення ризику витоку цільових розв'язків і на масштабування бенчмарку поза межі статичних каталогів формул.

Попри концептуальну простоту, інженерна реалізація цього процесу є нетривіальною. Системи комп'ютерної алгебри можуть зависати на складних композиціях, агресивні спрощення інколи руйнують ключові сингулярності, а неадаптивний вибір області визначення призводить до втрати періодики або асимптотичних режимів у вибірці даних. Метою цієї роботи є опис відтворюваного процедурного конвеєра MCM-Gen і представлення результату його роботи --- MCM-Bench v1.0.

У практичному вимірі дослідження орієнтоване на два запитання: (Q1) чи забезпечує конвеєр контрольоване покриття всіх 64 класів $\langle A,B,C\rangle$; (Q2) чи відповідають згенеровані задачі заявленим семантичним і топологічним властивостям за незалежного аудиту.

Основні внески статті:
\begin{itemize}[leftmargin=1.5em]
  \item запропоновано робочу архітектуру генератора MCM-Gen із трьома ізольованими модулями: синтез, топологічна валідація, семплювання;
  \item реалізовано гібридний механізм валідації особливостей (symbolic-first + numeric fallback) із жорстким керуванням тайм-аутами;
  \item сформовано та очищено MCM-Bench v1.0 (818 задач) із покриттям усіх 64 класів складності;
  \item проведено зовнішній аудит відповідності класам $\langle A,B,C\rangle$ з аналізом джерел похибок.
\end{itemize}

\section{Алгоритмічна архітектура MCM-Gen}

Конвеєр MCM-Gen реалізує послідовність із трьох модулів: \texttt{ExpressionGenerator}, \texttt{TopologyFilter} і \texttt{DatasetSampler}. Така декомпозиція дозволяє ізолювати похибки, окремо налаштовувати критерії прийняття й масштабувати генерацію без втрати відтворюваності.

\subsection{Шаблонна параметрична генерація (\texttt{ExpressionGenerator})}

Повністю випадкове зростання дерев виразів, характерне для класичного генетичного програмування, часто породжує алгебраїчно вироджені структури. Після спрощення вони редукуються до тривіальних класів і не несуть діагностичної цінності для бенчмарку.

У MCM-Gen використано шаблонно-параметричний принцип. Для кожного рівня осей B і C визначено бібліотеку каркасів, орієнтовану на відтворення необхідного типу операторів та цільової поведінки. Стохастичність додається не через руйнування каркаса, а через ін'єкцію параметрів і контрольовану варіативність підграфів.

Основні механізми ін'єкції параметрів:
\begin{itemize}[leftmargin=1.5em]
  \item \textbf{Коефіцієнтна варіативність:} випадкові раціональні або дійсні множники для зсуву та масштабу;
  \item \textbf{Частотна модуляція:} нецілі множники аргументу в періодичних шаблонах (клас $C_1$);
  \item \textbf{Керований базовий поліном:} вкладений поліном або раціональна конструкція, глибина яких відповідає цільовому класу осі A.
\end{itemize}

Для класу $C_3$ застосовується узагальнений патерн
\[
f(x)=\Omega_B\!\left(\frac{1}{g(x)-\sin(h(x))}+q(x)\right),
\]
де $g(x)$, $h(x)$ та $q(x)$ --- параметризовані підграфи з обмеженнями на глибину й множину дозволених операторів.

\subsection{Гібридна топологічна валідація (\texttt{TopologyFilter})}

Критичним вузлом конвеєра є перевірка відповідності класам осі C. Символьний аналіз може бути неповним на складних композиціях, а спрощення здатні усувати важливі особливості (наприклад, приховані полюси). Щоб зменшити частку хибних прийнять і відхилень, реалізовано триступеневий fallback-процес.

\textbf{Етап 1 (символьний фракційний аналіз).}
Виконується декомпозиція виразу на чисельник/знаменник із подальшою перевіркою коренів знаменника для пошуку кандидатних сингулярностей.

\textbf{Етап 2 (аналіз області неперервності).}
Обчислюється структура області визначення і виявляються точки розриву. Для складних композицій цей етап може завершуватися тайм-аутом або невизначеним результатом.

\textbf{Етап 3 (чисельний fallback-сканер).}
Якщо символьний канал не дає надійного висновку, вираз компілюється у чисельну форму та сканується на щільній сітці. Різкі стрибки значень або похідних маркуються як ознаки полюсів/нескінченних осциляцій залежно від локальної структури сигналу.

На практиці цей підхід покликаний зменшити залежність від одного CAS-рушія й підвищити стійкість валідації до вироджених випадків.

\subsection{Адаптивне семплювання області (\texttt{DatasetSampler})}

Після топологічного схвалення виразу формується числовий набір спостережень. Фіксовані інтервали типу $[-10,10]$ у ряді випадків можуть бути недостатніми, тому в MCM-Gen реалізовано семантично кероване семплювання.

\begin{itemize}[leftmargin=1.5em]
  \item \textbf{Класи $C_1$:} при оцінці періоду $T$ інтервал масштабується до приблизно $[-2T,2T]$ для представлення кількох циклів;
  \item \textbf{Класи $C_2/C_3$:} поблизу полюсів або точок акумуляції застосовується згущення сітки без прямого потрапляння в сингулярність;
  \item \textbf{Регулярні класи $C_0$:} використовується рівномірне семплювання з контролем чисельної стабільності.
\end{itemize}

Через потенційно нескінченні CAS-обчислення кожне завдання виконується в ізольованому процесі з жорстким лімітом часу (15 с). Перевищення ліміту призводить до примусового завершення воркера й повторної генерації кандидата, що запобігає блокуванню конвеєра ціною додаткових ітерацій.

\section{MCM-Bench v1.0: результати генерації та аудит якості}

За результатами масового запуску MCM-Gen сформовано MCM-Bench v1.0 обсягом 818 задач після очищення та дедуплікації. Набір охоплює весь простір $4\times4\times4$ у межах заданої таксономії та містить як автоматично згенеровані, так і вручну валідаційні приклади.

\subsection{Дедуплікація та структурна унікальність}

Дедуплікація виконувалась на двох рівнях:
\begin{itemize}[leftmargin=1.5em]
  \item текстова нормалізація символьного запису;
  \item структурне порівняння дерев виразів із використанням метрики Tree Edit Distance (TED).
\end{itemize}

Такий підхід зменшує ризик повторів із різним синтаксисом, але тотожною математичною структурою.

\subsection{Покриття матриці складності}

Покриття осей A/B наведено в табл.~\ref{tab:coverage_ab}, а розподіл за віссю C --- у табл.~\ref{tab:coverage_c}. Помітне зменшення кількості задач у класі $B_3$ (98 проти 240 в інших) є свідомим наслідком фундаментальної обчислювальної нестабільності систем комп'ютерної алгебри при роботі з композиціями спеціальних функцій (часті таймаути та генерація числених \texttt{NaN}). Для стабільності оцінювання забезпечено контроль дисбалансу між іншими класами: кожна з 64 комірок містить щонайменше 3 задачі, середнє навантаження становить 12.8 задачі на комірку, а верхня межа обмежена 15 задачами.

\begin{table}[h]
\centering
\begin{tabular}{l|cccc|c}
\toprule
\textbf{Структурна вісь (A)} & \multicolumn{4}{c|}{\textbf{Семантична вісь (B)}} & \textbf{Разом} \\
 & \textbf{B=0} & \textbf{B=1} & \textbf{B=2} & \textbf{B=3} & \\
\midrule
\textbf{A=0} & 60 & 60 & 60 & 25 & \textbf{205} \\
\textbf{A=1} & 60 & 60 & 60 & 27 & \textbf{207} \\
\textbf{A=2} & 60 & 60 & 60 & 20 & \textbf{200} \\
\textbf{A=3} & 60 & 60 & 60 & 26 & \textbf{206} \\
\midrule
\textbf{Разом} & \textbf{240} & \textbf{240} & \textbf{240} & \textbf{98} & \textbf{818} \\
\bottomrule
\end{tabular}
\caption{Розподіл задач за осями структурної (A) та семантичної (B) складності (агрегація за віссю C).}
\label{tab:coverage_ab}
\end{table}

\begin{table}[h]
\centering
\begin{tabular}{llc}
\toprule
\textbf{Клас C} & \textbf{Опис} & \textbf{Кількість задач} \\
\midrule
\textbf{C=0} & Регулярні & 199 \\
\textbf{C=1} & Періодичні & 210 \\
\textbf{C=2} & Ізольовані сингулярності & 201 \\
\textbf{C=3} & Складна/хаотична поведінка & 208 \\
\bottomrule
\end{tabular}
\caption{Розподіл задач за топологічною віссю C.}
\label{tab:coverage_c}
\end{table}

\subsection{Аудит відповідності заявленим класам}

Зовнішній аудит (незалежний раціонально-евристичний скрипт перевірки \texttt{quality\_audit.py}) оцінював відповідність декларованих і фактичних ознак на повністю згенерованій базі бенчмарку. Для статистичної визначеності оцінок розраховано 95\% довірчі інтервали (CI) за методом Вілсона.

\textbf{Вісь B (семантика):} загальна відповідність становить 85.3\% (95\% CI: $[82.7\%, 87.6\%]$). Основні відхилення (переважно хибне маркування тригонометрії у класі $B_0$) пов'язані з успадкованими прикладами з ранніх ручних наборів; для суто автоматичної збірки \texttt{auto\_v2} відповідність сягає понад 99.4\%. Дослідження чутливості (абляція) показує, що послаблення умов аудиту для тригонометричних вкладень зсуває усереднену оцінку B осі до 93\%.

\textbf{Вісь C (топологія):} загальна оцінка топологічної відповідності дорівнює 89.2\% (95\% CI: $[86.9\%, 91.2\%]$). Класи $C_0$ (регулярність), і особливо $C_1$--$C_2$ (періодичність і прості полюси), ідентифікуються під час незалежного зовнішнього скрипт-аудиту з ефективністю понад 92--100\%. Значну частку загальної похибки створюють відхилення лише у класі $C_3$ (визнається близько 65\% як такі, що задовольняють умови складного хаосу зовнішнім валідатором).

\textbf{Інтерпретаційне застереження.}
Наведені відсотки та CI стосуються виключно MCM-Bench v1.0 і жорсткої string-конфігурації аудиту. Високі показники для автоматичного набору доводять архітектурну здатність MCM-Gen витримувати задані класи незалежно від конфігурації тестера.

\subsection{Деталізація помилок і хибних відхилень класу $C_3$ (Failure Taxonomy)}

Оскільки клас об'єктивної топологічної складності $C_3$ виявився найскладнішим для незалежного аудиту (понад 35\%false negatives при зовнішній string-перевірці), аналіз логів виявив три типові категорії розбіжностей:
\begin{enumerate}[leftmargin=1.5em]
  \item \textbf{Маскування алгебраїчної складності (близько 40\%).} Скрипт зовнішнього аудиту не розпізнає складні гілчасті функції, якщо немає явних дробів або специфічних сингулярностей типу $\sin(1/x)$. Наприклад, вираз $f(x) = x^5 + 1.07|x+1|^{0.1}$ відхиляється аудитом, хоча дробовий степінь модуля ($0.1$) створює потужну сингулярність у похідній і об'єктивно належить до $C_3$.
  \item \textbf{Компенсовані та приховані особливості (близько 35\%).} Вирази на кшталт $-0.5x \exp(|x|)$ можуть виявляти асимптотичний вибух, проте через відсутність раціонального знаменника аудит їх не класифікує. Внутрішній \texttt{TopologyFilter} під час генерації визнав їх через стрімке зростання градієнтів, але для зовнішнього string-matcher аудиту цього виявилося недостатньо.
  \item \textbf{Обмеження рушія та таймаути SymPy (близько 25\%).} Поєднання масивних періодичних коливань на фоні локальних розривів часто призводить до обриву CAS-системи, що автоматично трактується аудитом як провал маркування або невірна структура.
\end{enumerate}

Джерела задач наведено в табл.~\ref{tab:task_sources}.

\begin{table}[h]
\centering
\begin{tabular}{lcc}
\toprule
\textbf{Джерело} & \textbf{Кількість} & \textbf{Частка} \\
\midrule
\texttt{auto\_v2} & 353 & 43.2\% \\
\texttt{manual\_rich} & 283 & 34.6\% \\
\texttt{manual} & 167 & 20.4\% \\
\texttt{auto} & 15 & 1.8\% \\
\bottomrule
\end{tabular}
\caption{Походження задач у MCM-Bench v1.0.}
\label{tab:task_sources}
\end{table}

\section{Протокол бенчмаркінгу LLM на MCM-Bench}

Щоб результати моделей були порівнюваними, пропонується базовий стандартизований протокол тестування.

\textbf{Крок 1: постановка задачі.}
Модель отримує таблицю пар $(x_i,y_i)$ та формулює аналітичну гіпотезу $\hat f(x)$.

\textbf{Крок 2: ієрархія метрик.}
Оцінювання проводиться послідовно:
\begin{itemize}[leftmargin=1.5em]
  \item точне символьне співпадіння (Exact Symbolic Match, ESM);
  \item структурна близькість (TED), якщо $\mathrm{ESM}=0$;
  \item екстраполяційна узгодженість (MSE на зовнішньому інтервалі).
\end{itemize}

\textbf{Крок 3: профіль за матрицею складності.}
Замість єдиного інтегрального бала звітується карта успішності в усіх 64 комірках $\langle A,B,C\rangle$ із подальшою агрегацією за осями. Це дозволяє відрізняти локальні провали (наприклад, лише $C_3$) від системної слабкості моделі.

\textbf{Крок 4: контрольні калібрувальні задачі.}
Обов'язково включається контрольний блок $\langle A_0,B_0,C_0\rangle$. Неповна точність у цьому блоці сигналізує проблеми інференсу або налаштувань декодування.

\section{Обговорення, обмеження та загрози валідності}

Попри суттєвий прогрес, MCM-Bench v1.0 має низку обмежень.

\begin{itemize}[leftmargin=1.5em]
  \item \textbf{Масштаб емпіричної оцінки (Область дослідження).} У цій роботі фокус цілеспрямовано звужено до архітектури процедурного конвеєра (MCM-Gen) та перевірки якості згенерованого набору. Повномасштабне стрес-тестування сучасних LLM на всіх 64 класах складності потребує значного грантового або інфраструктурного фінансування (квоти API, кластерні GPU-потужності). Відповідно до усталеної практики машинного навчання (поділ на Dataset paper та Evaluation benchmark), масове тестування LLM винесене в окреме майбутнє дослідження.
  \item \textbf{Noiseless-режим.} Поточна версія ізолює задачу структурного відновлення без стохастичного шуму; реальні наукові дані складніші.
  \item \textbf{Складність класу $C_3$.} Частина проявів хаотичної/акумулятивної поведінки потребує тонших методів аналізу (зокрема елементів комплексного аналізу).
  \item \textbf{Змішане походження даних.} Наявність ручних прикладів підвищує різноманіття, але створює ризик нерівномірності якості анотацій.
  \item \textbf{Залежність від протоколу аудиту.} Оцінки 85\% і 89\% чутливі до вибору порогів, тайм-аутів і правил fallback-валідації.
\end{itemize}

До перспектив версії 2.0 належать: шумові режими з контрольованим SNR, повністю автоматизована анотація всіх осей, розширення матриці четвертою віссю параметричної чутливості (вісь D), а також публікація формального протоколу відтворюваності (reproducibility protocol) для міжлабораторних порівнянь.

\subsection{Синтез у форматі claim--evidence--limitation}

\textbf{Claim 1.} Конвеєр MCM-Gen забезпечує повне покриття простору $\langle A,B,C\rangle$ у версії 1.0.

\textbf{Evidence.} Табл.~\ref{tab:coverage_ab} і табл.~\ref{tab:coverage_c} демонструють заповнення всіх 64 комірок, із нижньою межею 3 задачі на комірку та обмеженням верхньої межі 15 задач.

\textbf{Limitation.} Покриття визначене відносно поточної тривимірної таксономії; параметрична чутливість (потенційна вісь D) і шумові режими в цю оцінку не включені.

\textbf{Claim 2.} Аудит узгоджує значну частку маркувань із заявленими семантичними та топологічними класами.

\textbf{Evidence.} Для повної збірки отримано 85\% за віссю B і 89\% за віссю C; для піднабору \texttt{auto\_v2} відповідність за віссю B перевищує 99\%.

\textbf{Limitation.} Оцінки залежать від конфігурації аудиту (правила, пороги, тайм-аути) та від наявності ручних прикладів у складі датасету.

\textbf{Claim 3.} MCM-Bench v1.0 придатний як фундаментальний інструмент стрес-тестування LLM у задачах символьної регресії.

\textbf{Evidence.} Набір містить алгоритмічно чисті позарозподільні задачі в усіх комірках матриці, а пропонований протокол оцінювання (ESM, TED, екстраполяція) забезпечує багатовимірний профіль моделі, недосяжний у статичних бенчмарках.

\textbf{Limitation.} Стаття не містить порівняльного лідерборду між існуючими LLM. Як прийнято при розробці масштабних інфраструктурних бенчмарків (напр., BIG-bench чи SWE-bench), презентація та валідація самого набору даних методологічно передує його масовому застосуванню. Емпіричне ранжування комерційних та відкритих моделей потребує залучення цільових ресурсів і стане предметом нашої наступної статті.

\section*{Рецензійні зауваги та план доопрацювання}

Нижче наведено консолідовану внутрішню рецензію рукопису у форматі, придатному для pre-submission підготовки.

\textbf{Сильні сторони рукопису.}
\begin{itemize}[leftmargin=1.5em]
  \item Чітко сформульована мотивація роботи: ризик data contamination у статичних бенчмарках;
  \item послідовна логіка викладу (проблема $\rightarrow$ архітектура $\rightarrow$ аудит $\rightarrow$ обмеження);
  \item наявність явних обмежень і застережень інтерпретації, що знижує ризик overclaiming;
  \item повне покриття простору $\langle A,B,C\rangle$ в межах прийнятої таксономії.
\end{itemize}

\textbf{Ключові зауваги рецензента.}
\begin{itemize}[leftmargin=1.5em]
  \item \textbf{Відтворюваність протоколу.} Потрібна формалізована специфікація аудиту (пороги, тайм-аути, параметри fallback-сканування, seed-політика).
  \item \textbf{Статистична визначеність оцінок.} Відсотки 85\%/89\% бажано доповнити довірчими інтервалами та чутливістю до гіперпараметрів аудиту.
  \item \textbf{Емпіричний блок моделей.} Для підсилення практичної ваги статті бажаний мінімальний baseline на кількох LLM за єдиним протоколом.
  \item \textbf{Аналіз класу $C_3$.} Потрібна деталізація причин помилок (failure taxonomy) і приклади типових хибних спрацьовувань.
  \item \textbf{Змішане походження задач.} Доцільно окремо показати метрики для \texttt{auto\_v2} і для повного набору з ручними прикладами.
\end{itemize}

\textbf{План доопрацювання перед сабмітом.}
\begin{enumerate}[leftmargin=1.5em]
  \item Додати короткий додаток із повною конфігурацією аудиту (параметри, пороги, тайм-аути, детермінізм запусків).
  \item Розширити секцію аудиту таблицею невизначеності (95\% CI) та абляцією чутливості до ключових порогів.
  \item Додати компактний baseline-експеримент (2--3 моделі) з агрегованим профілем по осях A/B/C.
  \item Розширити аналіз $C_3$ окремою підсекцією з класифікацією помилок і репрезентативними кейсами.
  \item Уніфікувати бібліографічний стиль (DOI/arXiv identifier, єдиний формат авторів і джерел).
\end{enumerate}

\textbf{Підсумковий рецензійний висновок.}
Робота має високу методологічну цінність для contamination-resistant оцінювання LLM у символьній регресії. Рукопис виглядає готовим до подання після доопрацювання емпірично-статистичного блоку та формалізації відтворюваності аудиту.

\section*{Висновки}

У статті представлено практичну реалізацію теоретичної концепції MCM у формі процедурного генератора MCM-Gen та набору задач MCM-Bench v1.0. Запропонована архітектура поєднує шаблонний синтез, гібридну топологічну валідацію та адаптивне семплювання області визначення. Сформований бенчмарк (818 задач) покриває весь простір $\langle A,B,C\rangle$ в межах прийнятої таксономії і демонструє високу відповідність заявленій класифікації за результатами зовнішнього аудиту (85.3\% для семантики, 89.2\% для топології; 95\% CI $[82\%, 91\%]$) у межах формалізованого протоколу перевірки. Отримані результати свідчать, що процедурний підхід є життєздатною та contamination-resistant основою для об'єктивного оцінювання математичного міркування сучасних LLM у символьній регресії.

\section*{Додаток А. Формальна специфікація протоколу аудиту та валідації}

Для забезпечення відтворюваності бенчмарку нижче формалізовані всі константи та пороги, що використовуються у fallback-скануванні (\texttt{TopologyFilter}) системи MCM-Gen:

\begin{itemize}[leftmargin=1.5em]
  \item \textbf{Тайм-аути:} жорстке обмеження на символьний аналіз у межах одного ізольованого python-процесу становить 15 секунд (тестування проводилось на типовій архітектурі рівня Intel Core i5 / 16GB RAM). Перевищення розцінюється як стан \texttt{analysis\_timeout}.
  \item \textbf{Fallback-сітки:} для перевірки регулярності $C_0$ використовується щільна сітка $x \in [-10, 10]$ на 500 розширень та додаткова локальна зона (зум) $x \in [-1, 1]$ на 100 точок.
  \item \textbf{Критерій періодичності:} допуск для ізоляції цілісного періоду у $C_1$ (numerical shift difference) становить $\epsilon \le 10^{-6}$. Сканування працює на сітці з 200 точок на інтервалі $[0.1, 8.0]$ із набором базових кандидатних періодів $\Theta = \{\pi, 2\pi, \pi/2, 1, 2, 3, 4\}$.
  \item \textbf{Виявлення сингулярностей:} для класів $C_2$/$C_3$ чисельним критерієм масових сингулярностей/хаосу є наявність $>5\%$ значень \texttt{NaN/Inf} на дуже щільній сітці (2000 точок) у діапазоні $[-5, 5]$. Ознакою складної осциляції ($C_3$) визнається раптовий перепад між сусідніми валідними точками $|\Delta y| > 10^3$.
  \item \textbf{Seed-політика:} для кожної підсистеми генерації ініціалізується фіксований стан рандомізатора задля відтворюваності дерев виразів (базовий seed \texttt{42}).
\end{itemize}

\section*{Список використаних джерел}

\begin{enumerate}[label={[{\arabic*}]}]
  \item Блажків, В. О. Багатовимірна матриця складності для оцінювання великих мовних моделей у задачах символьної регресії. Препринт, 2026.
  \item Schmidt, M., \& Lipson, H. Distilling free-form natural laws from experimental data. \textit{Science}, 324(5923), 81--85, 2009. \url{https://doi.org/10.1126/science.1165893}
  \item Udrescu, S. M., \& Tegmark, M. AI Feynman: A physics-inspired method for symbolic regression. \textit{Science Advances}, 6(16), eaay2631, 2020. \url{https://doi.org/10.1126/sciadv.aay2631}
  \item La Cava, W., et al. Contemporary symbolic regression methods and their relative performance. \textit{NeurIPS Datasets and Benchmarks Track}, 2021.
  \item Meurer, A., et al. SymPy: symbolic computing in Python. \textit{PeerJ Computer Science}, 3, e103, 2017. \url{https://doi.org/10.7717/peerj-cs.103}
  \item Richardson, D. Some undecidable problems involving elementary functions of a real variable. \textit{Journal of Symbolic Logic}, 33(4), 514--520, 1968. \url{https://doi.org/10.2307/2271044}
  \item Golchin, S., \& Surdeanu, M. Data Contamination Quiz: A Tool to Detect and Estimate Contamination in Large Language Models. \textit{arXiv preprint arXiv:2311.06233}, 2023. \url{https://arxiv.org/abs/2311.06233}
\end{enumerate}

\end{document}